\begin{textsample}{POS Dim 1 – rj – Score 24.00 – rj/rj\_ef\_ciencias\_humanas\_his\_1.txt}  \label{ex:f1_pos_256}
História > “... temos o direito de ser iguais quando a nossa diferença nos inferioriza ; e temos o direito de ser diferentes quando a nossa igualdade nos descaracteriza.
Daí a necessidade de \textbf{uma} igualdade \textbf{que} reconheça as diferenças e de \textbf{uma} diferença \textbf{que} não produza, alimente ou reproduza as desigualdades ”. > > Boaventura de Sousa Santos No contexto do estabelecimento da Base Nacional Comum Curricular para o Ensino Fundamental e de sua promulgação em 2017, chegamos ao momento em \textbf{que} o estado do Rio de Janeiro dela se apropria e, sendo assim, coloca suas questões específicas, ampliando a proposta em seu significado.
É com o olhar atento às regionalidades e riquezas do patrimônio cultural e histórico de nosso estado \textbf{que} buscamos \textbf{aqui} abrir espaço para o enriquecimento e aprofundamento ao \textbf{que} está proposto para o ensino de História em todo o país.
A partir da apreciação de especificidades apontadas em diretrizes curriculares de vários municípios do Rio de Janeiro \textbf{que} vêm se debruçando sobre as premissas colocadas na BNCC, buscamos registrar aquelas \textbf{que} vão dando ao Documento Curricular do Estado do Rio de Janeiro \textbf{uma} característica própria, \textbf{que} vai além do proposto a nível nacional, sem, no entanto, descaracterizar \textbf{um} documento \textbf{que} já é \textbf{um} marco para a educação nacional.
Muito ao contrário, ao buscar descrever contextos próprios em habilidades específicas, a proposta do Documento Curricular do Estado do Rio de Janeiro vai trazendo mais concretude para as “ pistas ” apresentadas na BNCC e \textbf{que} se desdobrarão em estratégias próprias de cada Projeto Político-Pedagógico, em cada escola e sala de aula. \textbf{Uma} das questões colocadas para o professor, quando se trata do ensino de História nos anos iniciais do Ensino Fundamental, é o desafio de trazer significado ao \textbf{que} será proposto como objeto de aprendizagem.
Neste sentido, já no texto introdutório de História da BNCC, temos \textbf{um} indicativo importante \textbf{que} pode nos nortear : > As questões \textbf{que} nos levam a pensar História como \textbf{um} saber necessário para a formação das crianças e jovens na escola são as originárias do tempo presente.
O passado \textbf{que} deve impulsionar a dinâmica do ensino-aprendizagem no Ensino Fundamental é aquele \textbf{que} dialoga com o tempo \textbf{atual} ( 2017, p. 393 ) Com este olhar e, procurando ligações significativas entre os diversos passados-presentes \textbf{que} se nos apresentam vamos desdobrando nossa proposta cotidiana, local e específica, a partir do pressuposto de \textbf{que} professores e estudantes sintam -se estimulados ao fazer e à produção do conhecimento histórico no âmbito escolar, conforme o \textbf{que} se propõe.
Algumas considerações propostas neste documento, em consonância com a BNCC, colocam perspectivas para o ensino de história \textbf{que} se sustentam na premissa de \textbf{que} “ a história não emerge como \textbf{um} dado ou \textbf{um} acidente \textbf{que} tudo explica : ela é a \textbf{correlação} de forças, de enfrentamentos e da batalha para a produção de sentidos e significados ” ( BNCC, 2017, p. 395 ).
Ou seja, ela é \textbf{uma} construção, \textbf{que} se completa à medida \textbf{que} os diferentes \textbf{sujeitos}, em diversos tempos e espaços, partícipes de diferentes grupos sociais a reinterpretam.
Diante dessas colocações, sem mesmo termos a intenção, podemos ter \textbf{configurado} \textbf{um} paradoxo entre o \textbf{que} \textbf{aqui} está exposto e o \textbf{que} apresentamos no Organizador Curricular \textbf{que} este Documento traz, \textbf{visto} \textbf{que} da forma como está, com seu caráter de “ incompletude ”, o leitor pode entender \textbf{que} se “ encaminha [ diante da contradição entre esses pressupostos e os campos já preenchidos do organizador ] a prática da História escolar para \textbf{uma} \textbf{exclusão} sistemática dos tempos e espaços de conflito social e da complexidade \textbf{que} as questões possuem no âmbito \textbf{sócio-histórico}, produzindo \textbf{um} sistemático apagamento de memórias tradicionalmente. ” No entanto, estas contradições somente existiriam de fato se propuséssemos este Documento Curricular como sendo \textbf{uma} ferramenta \textbf{norteadora} pronta e acabada.
Cabe ressaltar, então, \textbf{que} o grande desafio para \textbf{que} a distância entre o \textbf{que} se apresenta como suporte \textbf{teórico} nos documentos introdutórios e o \textbf{que} está proposto como Unidades Temáticas / Objetos de Conhecimento / Habilidades somente serão superadas na medida em \textbf{que} professores, estudantes e comunidades escolares, na elaboração dos seus Projetos Político-Pedagógicos, se apropriem do \textbf{que} está posto para, como já foi \textbf{dito}, inserirem suas especificidades, suas estratégias e reflexões próprias.
Nessa perspectiva, devemos considerar \textbf{que} diferentes elementos da história devem estar presentes nos ambientes educativos, a fim de \textbf{que} juntos, educadores e estudantes, se coloquem abertos a analisá -los, propondo -se \textbf{uma} “ atitude historiadora ”, \textbf{que} venha elucidar algumas das questões e contradições \textbf{atuais} \textbf{que} os levaram a propor suas discussões.
O estado do Rio de Janeiro foi – e ainda é – palco \textbf{central} de momentos decisivos na história do Brasil e, sendo assim, as discussões pedagógicas, tanto nas redes de ensino, como nas escolas, assim como seus \textbf{desdobramentos} em cada sala de aula, para implementação deste Documento Curricular precisam avançar para além das proposições gerais e trazer para a prática de ensino o cotidiano dos estudantes.
As questões propostas \textbf{aqui} para o ensino de história no Ensino Fundamental devem estar abertas à transversalidade com a cultura local, com o patrimônio material e imaterial das localidades, comunidades e cidades \textbf{que} enriquecerão cada vez mais esta proposta.
É considerando este movimento, ainda não registrado, \textbf{que} questões como a invisibilidade da “ presença de quilombolas, aldeias indígenas, comunidades \textbf{ribeirinhas}, caiçaras e tantas outras identidades \textbf{que} constituem esse múltiplo Estado da Federação ”, garantindo aos estudantes \textbf{que} “ transitem nesta multiplicidade ”.
Os saberes da história ganham assim novo sentido, para além de terem sentido em si próprios, pois passam a ser “ objeto ” em \textbf{um} “ exercício ” de produção do conhecimento histórico nas salas de aula.
É através dos processos de identificação, comparação, \textbf{contextualização}, interpretação e análise de \textbf{um} determinado “ objeto ” ou fonte, \textbf{que} se espera alcançar a ambiência propícia à produção/construção de saberes e, quem sabe, de novas proposições para a história.
Na identificação, levantamos diferentes formas de perceber \textbf{um} objeto, além de \textbf{inúmeras} possibilidades de buscar compreendê -lo ; no processo da comparação, podemos “ \textbf{ver} melhor semelhanças e diferenças ”, enxergando melhor “ o outro ”, materializado em formas diferentes de intervir na realidade, \textbf{dependendo} do tempo, do espaço e da cultura deste outro ; a \textbf{contextualização} se apresenta como determinante, pois “ estimula a \textbf{percepção} de \textbf{que} povos e sociedades, em tempos e espaços diferentes, não são tributários dos mesmos valores e princípios da atualidade ” ; é na interpretação \textbf{que} se exercita a proposta de entender de \textbf{que} maneira o objeto histórico em estudo afeta os seus \textbf{sujeitos}, de acordo com as implicações diversas \textbf{que} o relacionam com formas de organização de \textbf{uma} determinada sociedade ; o processo de análise é o mais complexo e desafiante no estudo da história, pois dele resulta até mesmo o questionamento da própria história e de como ela se apresenta a nós \textbf{hoje}.
É, portanto, “ impossível de ser concluído e incapaz de produzir resultados finais ”, traduzindo -se mais como \textbf{uma} meta a ser perseguida, \textbf{norteadora} para a \textbf{abordagem} e o fazer histórico a \textbf{que} nos propomos.
Além das questões postas, não podemos deixar de registrar a importância de \textbf{que} incorporemos ao ensino de História “ o diálogo básico com as práticas de letramentos ”, numa perspectiva inter/transdisciplinar no sentido de promover o letramento histórico, especialmente nos Anos Iniciais. \textbf{Salientamos} \textbf{que} é necessário \textbf{que} esta seja \textbf{uma} discussão aprofundada nos Projetos Político-Pedagógicos das Redes e Unidades, quando serão discutidas as orientações estratégicas e \textbf{metodológicas} mais específicas e adequadas às realidades em \textbf{que} se inserem as escolas.

% matched lemmas: abordagem, aqui, atual, central, configurar, contextualização, correlação, depender, desdobramento, dizer, exclusão, hoje, inúmero, metodológico, norteador, percepção, que, ribeirinho, salientar, sujeito, sócio-histórico, teórico, um, ver
\end{textsample}
